\documentclass[11pt]{article}
\usepackage[margin=1in]{geometry}
\usepackage{amsmath,amssymb,amsthm}
\usepackage{graphicx}
\usepackage[hidelinks]{hyperref}
\usepackage{xcolor}
\usepackage{microtype}

\newtheorem{theorem}{Theorem}
\newcommand{\T}{\mathbb T}

\title{A Sparse Non-Unimodular Equality Case\\
for the Complex Spencer Problem in Dimension Four}
\author{Autonomous discovery packet for arXiv:2602.12868}
\date{13 August 2026}

\begin{document}
\maketitle

\begin{center}
\fbox{\parbox{0.91\textwidth}{
\textbf{Verdict.} Candidate full solution, likely valid, to the source's
explicit equality-case question.  This note constructs non-unimodular rows
attaining the sharp value $\sqrt4=2$.  It does \emph{not} prove the universal
$n=4$ instance of Conjecture~1.2.
}}
\end{center}

\section{Source question}

Tomasz Kobos and Marin Varivoda, \emph{A Complex Analogue of Spencer's Six
Standard Deviations Theorem and the Complex Banach--Mazur Distance},
arXiv:2602.12868 (2026), formulate Conjecture~1.2: if
$a_1,\ldots,a_n\in\mathbb C^n$ satisfy $\|a_i\|_\infty\leq1$, then some
$x\in\T^n$ satisfies $|\langle x,a_i\rangle|\leq\sqrt n$ for all $i$.
They prove it for $n=2,3$.  At the end of page~26 they write:

\begin{quote}
``We do not know whether, as in the case $n=3$, there also exist
non-unimodular vectors achieving equality.''
\end{quote}

Here ``unimodular'' means that every coordinate has modulus one, and the source
uses $\langle x,y\rangle=\sum_jx_j\overline{y_j}$.  The complete source context
is reproduced below.

\begin{center}
\includegraphics[width=0.98\textwidth]{figures/open_problem_crop.png}
\end{center}

\section{Result}

For a row system $A=(a_i)_{i=1}^4$, put
\[
 D(A):=\min_{x\in\T^4}\max_{1\leq i\leq4}|\langle x,a_i\rangle|.
\]
The source calls $A$ an equality case in dimension four when $D(A)=2$.

\subsection*{Proof intuition}

Take the four sign patterns of two variables and append a zero coordinate.
The resulting correlations are the four vertices $1\pm z\pm w$ of a planar
parallelogram translated by $1$.  If all four vertices entered the open disk of
radius two, elementary expansion would force the three directions
$1,z,w$ to have pairwise line angles strictly larger than $60^\circ$.  Three
lines cannot be packed that way into a projective circle of circumference
$180^\circ$.  Cube-root directions attain the boundary exactly.

\begin{theorem}[Non-unimodular $n=4$ extremizer]
For $(\varepsilon,\delta)\in\{\pm1\}^2$, let
\[
 a_{\varepsilon,\delta}=(1,\varepsilon,\delta,0)\in\mathbb C^4.
\]
Then
\[
 \min_{x\in\T^4}\max_{\varepsilon,\delta\in\{\pm1\}}
 |\langle x,a_{\varepsilon,\delta}\rangle|=2.
\]
Consequently, dimension four admits a non-unimodular equality case for
Conjecture~1.2.
\end{theorem}

\begin{proof}
Let $x=(x_1,x_2,x_3,x_4)\in\T^4$.  Multiplying all correlations by
$\overline{x_1}$ does not change their moduli, so set
$z=x_2\overline{x_1}$ and $w=x_3\overline{x_1}$.  Thus $|z|=|w|=1$ and
\[
 |\langle x,a_{\varepsilon,\delta}\rangle|
   =|1+\varepsilon z+\delta w|.
\]
Write
\[
 a=\operatorname{Re}z,\qquad b=\operatorname{Re}w,\qquad
 c=\operatorname{Re}(z\overline w).
\]
Expansion gives
\begin{equation}\label{eq:expand}
 |1+\varepsilon z+\delta w|^2
 =3+2q_{\varepsilon,\delta},\qquad
 q_{\varepsilon,\delta}=\varepsilon a+\delta b+\varepsilon\delta c.
\end{equation}

Suppose, for contradiction, that all four moduli were strictly less than $2$.
Then every $q_{\varepsilon,\delta}<1/2$.  Adding the two inequalities with a
fixed first sign gives
\[
 q_{+, +}+q_{+,-}=2a<1,
 \qquad q_{-,+}+q_{-,-}=-2a<1,
\]
and hence $|a|<1/2$.  Fixing the second sign in the same way yields
$|b|<1/2$.  Finally,
\[
 q_{+,+}+q_{-,-}=2c<1,
 \qquad q_{+,-}+q_{-,+}=-2c<1,
\]
so $|c|<1/2$.

Regard $1,z,w$ as unit vectors in the real plane.  Their three pairwise real
inner products are $a,b,c$.  Passing to unoriented lines places their arguments
on the projective circle $\mathbb R/\pi\mathbb Z$.  Three points divide this
circle of circumference $\pi$ into three gaps, one of which is at most
$\pi/3$.  The corresponding two lines therefore have acute angle at most
$\pi/3$, so the absolute value of their inner product is at least
$\cos(\pi/3)=1/2$.  This contradicts $|a|,|b|,|c|<1/2$.  Hence every
$x\in\T^4$ has at least one correlation of modulus at least $2$.

For the reverse inequality let $\omega=e^{2\pi i/3}$ and take
$x=(1,\omega,\omega^2,1)$.  Since $1+\omega+\omega^2=0$, the four correlation
moduli are
\[
 |1+\omega+\omega^2|,\quad |1+\omega-\omega^2|,\quad
 |1-\omega+\omega^2|,\quad |1-\omega-\omega^2|
   =0,2,2,2.
\]
Thus $D(A)\leq2$, completing the proof.  Each row has a zero fourth coordinate,
so none is unimodular.
\end{proof}

\section{Verification}

The proof is exact and computation-free.  The accompanying script
\texttt{code/verify\_extremizer.py} independently checks:
\begin{itemize}
\item the four exact cube-root witness values $0,2,2,2$;
\item one million uniformly sampled phase pairs, all satisfying the lower
bound to floating-point tolerance; and
\item a differential-evolution minimization returning $2$ to better than
$10^{-10}$.
\end{itemize}
The key identity \eqref{eq:expand} can also be checked directly term by term.
No numerical statement is used in the proof.

\section{Scope, novelty, and review recommendation}

This theorem fully answers the narrow existence question quoted above.  It
does not classify equality cases, prove Conjecture~1.2 for arbitrary four-row
systems, or settle either Banach--Mazur conjecture in the source.

On 13 August 2026, bounded searches used arXiv:2602.12868, the exact paper
title, the exact phrase ``non-unimodular vectors achieving equality,'' and
combinations of ``complex Spencer,'' ``dimension 4,'' and ``non-unimodular
equality.''  They found the source preprint and mirrors but no later resolution,
follow-up, or erratum.  This supports, but does not certify, novelty.

Human review should focus on (i) the strict-inequality step from four open
radius-two disks to $|a|,|b|,|c|<1/2$, (ii) agreement with the source inner
product convention, and (iii) whether the sparse Walsh configuration is known
under another formulation.  Subject to those checks, this is ready to report
to the authors as a short affirmative answer to their equality-case question.

\begin{thebibliography}{1}
\bibitem{KV}
T.~Kobos and M.~Varivoda,
\emph{A Complex Analogue of Spencer's Six Standard Deviations Theorem and the
Complex Banach--Mazur Distance}, arXiv:2602.12868, 2026.
\end{thebibliography}

\end{document}
